\chapter{Logit e Probit}
\label{cap:logit}
% ── índice ──────────────────────────────────────────────────────
\index{logit|textbf}
\index{probit|textbf}
\index{logit@\texttt{logit()}|textbf}
\index{probit@\texttt{probit()}|textbf}
\index{máxima verossimilhança|textbf}
\index{efeitos marginais|textbf}
\index{margins@\texttt{margins()}|textbf}
\index{odds ratio}

% ─────────────────────────────────────────────────────────────────────────────
\section{O modelo de variável latente}

Quando a variável dependente é binária ($y \in \{0,1\}$), o modelo linear
produz previsões fora do intervalo $[0,1]$ e erros heterocedásticos por
construção. A solução é modelar a probabilidade de $y=1$ via uma função
de distribuição acumulada:

\[
  P(y_i = 1 \mid x_i) = F(x_i'\beta)
\]

A motivação vem do modelo de variável latente:
\[
  y_i^* = x_i'\beta + \varepsilon_i, \qquad y_i = \mathbf{1}[y_i^* > 0]
\]

A distribuição de $\varepsilon$ determina o modelo:

\begin{center}
\begin{tabular}{lll}
\toprule
\textbf{Modelo} & \textbf{Distribuição de $\varepsilon$} & \textbf{$F(\cdot)$} \\
\midrule
Logit  & Logística$(0, \pi^2/3)$ & $\Lambda(z) = \dfrac{1}{1 + e^{-z}}$ \\[8pt]
Probit & Normal$(0, 1)$           & $\Phi(z) = \displaystyle\int_{-\infty}^{z} \phi(t)\,dt$ \\
\bottomrule
\end{tabular}
\end{center}

A estimação é por Máxima Verossimilhança (MLE). A log-verossimilhança é:
\[
  \ell(\beta) = \sum_{i=1}^{n}
    \left[y_i \ln F(x_i'\beta) + (1-y_i)\ln\bigl(1 - F(x_i'\beta)\bigr)\right]
\]

% ─────────────────────────────────────────────────────────────────────────────
\section{Interpretação dos coeficientes}

Os coeficientes do Logit e Probit \textbf{não são efeitos marginais} — são
parâmetros da variável latente, sem unidade interpretável diretamente.
Para $x_j$ contínuo, o efeito marginal sobre a probabilidade é:

\[
  \frac{\partial P(y=1 \mid x)}{\partial x_j}
  = f(x'\beta)\,\beta_j
\]

onde $f = F'$ é a densidade da distribuição escolhida. Como $f(x'\beta)$
varia entre observações, reporta-se o \textbf{Efeito Marginal Médio} (AME):

\[
  \mathrm{AME}_j = \frac{1}{n}\sum_{i=1}^{n} f(x_i'\beta)\,\beta_j
\]

\begin{nota}
  A relação entre os coeficientes de Logit e Probit é aproximadamente:
  \[
    \hat\beta_{\mathrm{Logit}} \approx \frac{\pi}{\sqrt{3}}\,\hat\beta_{\mathrm{Probit}}
    \approx 1{,}81\,\hat\beta_{\mathrm{Probit}}
  \]
  Os AMEs dos dois modelos são tipicamente muito próximos, mesmo com
  coeficientes em escalas diferentes.
\end{nota}

% ─────────────────────────────────────────────────────────────────────────────
\section{Verificação por DGP}

\subsection*{Consistência do MLE}

Ao contrário dos modelos lineares dos capítulos anteriores, Logit e Probit
\emph{não admitem recuperação exata} dos parâmetros. A variável binária $y$
é gerada por amostragem de Bernoulli — há aleatoriedade irredutível mesmo
sem adicionar ruído extra. O que se verifica aqui é \textbf{consistência}:
com $n=1000$ o MLE deve convergir próximo dos parâmetros verdadeiros.

O DGP usa a parametrização logística com $\beta_0 = -1{,}0$ e $\beta_1 = 2{,}0$:
\[
  x_i \sim \mathcal{N}(0,1), \qquad
  p_i = \Lambda(-1{,}0 + 2{,}0\,x_i), \qquad
  y_i \sim \mathrm{Bernoulli}(p_i)
\]

Bernoulli individual é gerado por $y_i = \mathbf{1}[u_i < p_i]$, com
$u_i \sim \mathrm{Uniforme}(0,1)$:

\begin{lstlisting}[caption={DGP logístico — Logit e Probit simultâneos}]
seed(42)
input df
id
1
end
for i in 1..=10 { df = append(df, df) }
df |> filter(_n <= 1000)
generate df x  = rnormal()
generate df lp = -1.0 + 2.0*x
generate df p  = 1 / (1 + exp(-lp))
generate df u  = runiform()
generate df y  = u < p
let m_l = logit(y ~ x, df)
let m_p = probit(y ~ x, df)
esttab(m_l, m_p)
\end{lstlisting}

Os verdadeiros parâmetros são $(\beta_0, \beta_1) = (-1{,}0;\; 2{,}0)$ para
o DGP Logit com $x_i \sim \mathcal{N}(0,1)$. As estimativas Probit estão na
escala normal do índice latente; portanto, não devem ser numericamente idênticas
aos coeficientes Logit mesmo quando os dois modelos descrevem a mesma relação
qualitativa.

\subsection*{Efeitos marginais médios}

Os coeficientes são apenas parâmetros da variável latente. Para
interpretação econômica, usa-se \lstinline{margins}:

\begin{lstlisting}[caption={AME do Logit}]
margins(m_l, df)
\end{lstlisting}

O efeito marginal é positivo e varia com a probabilidade ajustada; ele é
reportado por \lstinline{margins(m_l, df)} na escala de probabilidade.

% ─────────────────────────────────────────────────────────────────────────────
\section{Sintaxe}

\begin{lstlisting}
logit(y ~ x1 + x2, df)
probit(y ~ x1 + x2, df)
margins(modelo, df)
\end{lstlisting}

\begin{lstlisting}[caption={Fluxo típico — escolha discreta}]
load "dados.csv" as df

let m_l = logit(empregado ~ educ + exp + sexo, df)
let m_p = probit(empregado ~ educ + exp + sexo, df)

esttab(m_l, m_p)
margins(m_l, df)
\end{lstlisting}

\section{Previsão de probabilidades}

\begin{lstlisting}
let m = logit(y ~ x, df)
let phat = predict(m, "pr")     // P(y=1|x) para cada obs
let yhat = predict(m, "class")  // classe prevista (0 ou 1)
\end{lstlisting}

\begin{nota}
  Use \lstinline{predict(m, "class")} com limiar padrão de 0{,}5.
  Em dados desbalanceados (prevalência baixa), ajuste o limiar com
  \lstinline{predict(m, "class", threshold=0.3)}.
\end{nota}
